\documentclass[10pt]{beamer}

% --- Modern Theme ---
\usetheme{metropolis}
\usecolortheme{default}
\setbeamercolor{background canvas}{bg=white}

% --- Packages ---
\usepackage[utf8]{inputenc}
\usepackage{booktabs}
\usepackage{array}
\usepackage{tabularx}
\usepackage{amsmath}
\usepackage{amssymb}
\usepackage[scale=2]{ccicons}
\usepackage{pgfplots}
\usepgfplotslibrary{dateplot}
\usepackage{xspace}

% --- Title Information ---
\title{Foundations of Artificial Intelligence}
\subtitle{Part VII: Introduction to Robotics \& Computer Vision (Overview)}
\institute{Department of Computer Science}
\date{}

\begin{document}

\maketitle

\begin{frame}{Lecture Roadmap}
    \setbeamertemplate{section in toc}[sections numbered]
    \tableofcontents[hideallsubsections]
\end{frame}

\begin{frame}{Learning Objectives}
    By the end of this lecture, you should be able to:
    \begin{itemize}
        \item Explain how robotics combines perception, reasoning, and action.
        \item Identify common sensors used by intelligent robots and autonomous systems.
        \item Describe the basic idea of image classification in computer vision.
        \item Discuss how AI enables autonomy in self-driving and other robotic platforms.
        \item Recognize practical limitations: uncertainty, safety, and real-time constraints.
    \end{itemize}
\end{frame}

\section{Robotics and Vision Overview}

\begin{frame}{What Is Robotics in AI?}
    \textbf{Robotics:} Building agents that sense, reason, and act in the physical world.

    \vspace{0.2cm}
    Core components:
    \begin{itemize}
        \item \textbf{Perception:} Understand environment from sensors.
        \item \textbf{Decision-making:} Plan and choose actions.
        \item \textbf{Control:} Execute actions through motors/actuators.
    \end{itemize}

    \vspace{0.2cm}
    \textbf{Computer Vision} is a key part of robotic perception:
    \begin{itemize}
        \item Detect lanes, objects, humans, obstacles.
        \item Estimate scene structure and motion.
        \item Support safe navigation and interaction.
    \end{itemize}
\end{frame}

\begin{frame}{Why Robotics + Computer Vision Matters}
    Real systems where AI and CV are essential:
    \begin{itemize}
        \item Autonomous driving and driver-assistance.
        \item Warehouse and logistics robots.
        \item Surgical and assistive medical robots.
        \item Drones for mapping, inspection, and delivery.
        \item Smart manufacturing and quality control.
    \end{itemize}

    \vspace{0.2cm}
    \textbf{High-level loop:}
    \[
      \text{Sense} \rightarrow \text{Understand} \rightarrow \text{Plan} \rightarrow \text{Act} \rightarrow \text{Repeat}
    \]
\end{frame}

\begin{frame}{A Canonical Robotics Pipeline}
    \small
    \setlength{\tabcolsep}{8pt}
    \renewcommand{\arraystretch}{1.2}
    \begin{tabular}{>{\raggedright\arraybackslash}p{2.8cm} >{\raggedright\arraybackslash}p{7.2cm}}
        \toprule
        \textbf{Stage} & \textbf{Typical Tasks} \\
        \midrule
        Sensing & Capture data using camera, LiDAR, radar, IMU, GPS, ultrasound. \\
        Perception & Object detection, localization, scene understanding, tracking. \\
        Mapping & Build map or occupancy grid from observations. \\
        Planning & Path planning and behavior selection under constraints. \\
        Control & Convert plans into steering, velocity, and actuator commands. \\
        Feedback & Re-evaluate after each action to handle uncertainty. \\
        \bottomrule
    \end{tabular}
\end{frame}

\section{Perception}

\begin{frame}{Perception: From Raw Data to Meaning}
    \textbf{Perception} transforms sensor signals into useful world models.

    \vspace{0.2cm}
    Typical outputs:
    \begin{itemize}
        \item \textbf{Detection:} What objects are present?
        \item \textbf{Localization:} Where are they?
        \item \textbf{Classification:} What category does each object belong to?
        \item \textbf{Tracking:} How are objects moving over time?
    \end{itemize}

    \vspace{0.2cm}
    \textbf{Challenge:} Real-world data is noisy, incomplete, and ambiguous.
\end{frame}

\begin{frame}{Levels of Visual Perception}
    \begin{itemize}
        \item \textbf{Low-level vision:} edges, corners, gradients, keypoints.
        \item \textbf{Mid-level vision:} regions, contours, depth cues, motion cues.
        \item \textbf{High-level vision:} object categories, scene semantics, intent.
    \end{itemize}

    \vspace{0.2cm}
    Classical examples:
    \begin{itemize}
        \item Stereo vision: estimate depth from disparity.
        \item Optical flow: estimate motion from frame-to-frame changes.
    \end{itemize}

    \vspace{0.2cm}
    Modern systems combine classical geometry with learned models.
\end{frame}

\begin{frame}{Perception Under Uncertainty}
    Why uncertainty appears:
    \begin{itemize}
        \item Sensor noise and calibration errors.
        \item Occlusion (objects partially hidden).
        \item Changing illumination and weather.
        \item Dynamic scenes with many moving agents.
    \end{itemize}

    \vspace{0.2cm}
    Common strategies:
    \begin{itemize}
        \item Probabilistic filtering (e.g., Kalman-like updates).
        \item Multi-sensor fusion for robustness.
        \item Temporal smoothing and confidence estimates.
    \end{itemize}
\end{frame}

\section{Sensors}

\begin{frame}{Sensors in Intelligent Robots}
    \textbf{Sensors} provide observations used for state estimation and planning.

    \vspace{0.2cm}
    Key categories:
    \begin{itemize}
        \item \textbf{Vision:} RGB/mono/stereo cameras.
        \item \textbf{Range:} LiDAR, sonar, infrared depth sensors.
        \item \textbf{Motion:} IMU (accelerometer + gyroscope), wheel encoders.
        \item \textbf{Global position:} GPS/GNSS.
        \item \textbf{Active safety sensing:} radar for distance and velocity.
    \end{itemize}
\end{frame}

\begin{frame}{Sensor Trade-offs}
    \small
    \setlength{\tabcolsep}{7pt}
    \renewcommand{\arraystretch}{1.2}
    \begin{tabularx}{\textwidth}{l l X}
        \toprule
        \textbf{Sensor} & \textbf{Strength} & \textbf{Limitation} \\
        \midrule
        Camera & Rich semantic detail, low cost & Sensitive to lighting and weather \\
        LiDAR & Accurate 3D distance measurements & Expensive, lower semantic detail \\
        Radar & Works in fog/rain, good velocity estimation & Lower spatial resolution \\
        IMU & Fast motion updates & Drift over long duration \\
        GPS & Global localization outdoors & Weak/absent indoors or urban canyons \\
        \bottomrule
    \end{tabularx}

    \vspace{0.2cm}
    \textbf{Takeaway:} No single sensor is enough; practical systems fuse multiple sensors.
\end{frame}

\begin{frame}{Sensor Fusion (Basic Idea)}
    \textbf{Goal:} Combine complementary sensor information to improve reliability.

    \vspace{0.2cm}
    Example in autonomous driving:
    \begin{itemize}
        \item Camera: recognizes traffic signs and lane markings.
        \item LiDAR: provides geometric obstacle boundaries.
        \item Radar: estimates relative speed of nearby vehicles.
        \item GPS + IMU: provides global position and smooth local motion estimate.
    \end{itemize}

    \vspace{0.2cm}
    Fusion often improves robustness when one modality fails.
\end{frame}

\section{Image Classification (Basic Idea)}

\begin{frame}{What Is Image Classification?}
    \textbf{Task:} Assign an image to one label from a predefined set.

    \vspace{0.2cm}
    Example classes: \{cat, dog, car, pedestrian, stop-sign\}.

    \vspace{0.2cm}
    Pipeline (simplified):
    \begin{enumerate}
        \item Collect labeled images.
        \item Preprocess (resize, normalize, augment).
        \item Train model to map pixels $\to$ class probabilities.
        \item Predict class with highest probability.
    \end{enumerate}
\end{frame}

\begin{frame}{From Features to Deep Learning}
    \textbf{Classical approach:}
    \begin{itemize}
        \item Hand-crafted features (edges, texture descriptors).
        \item Classifier (e.g., SVM, logistic regression).
    \end{itemize}

    \vspace{0.2cm}
    \textbf{Modern approach: CNNs}
    \begin{itemize}
        \item Learn hierarchical features directly from data.
        \item Early layers: edges and color blobs.
        \item Deeper layers: parts and object-level concepts.
    \end{itemize}

    \vspace{0.2cm}
    \textbf{Why CNNs work well:} local receptive fields + shared weights + translation tolerance.
\end{frame}

\begin{frame}{Classification Output and Decision}
    For image $x$, model outputs class probabilities:
    \[
      p(y=k\mid x), \quad k=1,\dots,K
    \]

    \vspace{0.2cm}
    Predicted label:
    \[
      \hat{y} = \arg\max_k p(y=k\mid x)
    \]

    \vspace{0.2cm}
    \textbf{Confidence matters:}
    \begin{itemize}
        \item High confidence: model is certain.
        \item Low confidence: defer to human or safer fallback policy.
    \end{itemize}

    \vspace{0.2cm}
    In safety-critical systems, decision policies include uncertainty thresholds.
\end{frame}

\begin{frame}{Beyond Classification}
    In robotics, classification alone is not enough.

    \vspace{0.2cm}
    Additional vision tasks:
    \begin{itemize}
        \item \textbf{Object detection:} classify + localize with bounding boxes.
        \item \textbf{Semantic segmentation:} classify each pixel.
        \item \textbf{Depth estimation:} infer distance from images.
        \item \textbf{Pose estimation:} infer orientation and key points.
    \end{itemize}

    \vspace{0.2cm}
    These tasks support navigation, manipulation, and human-robot interaction.
\end{frame}

\section{AI in Autonomous Systems}

\begin{frame}{AI Stack in Autonomous Systems}
    A practical autonomy stack often includes:
    \begin{itemize}
        \item \textbf{Perception models:} detect lanes, vehicles, pedestrians.
        \item \textbf{Prediction models:} forecast other agents' motion.
        \item \textbf{Planning models:} choose safe and efficient trajectories.
        \item \textbf{Control algorithms:} track trajectory under dynamics constraints.
        \item \textbf{Monitoring/safety layer:} override unsafe actions.
    \end{itemize}

    \vspace{0.2cm}
    System-level reliability matters more than any single model's accuracy.
\end{frame}

\begin{frame}{Example: Autonomous Driving Loop}
    \small
    \setlength{\tabcolsep}{7pt}
    \renewcommand{\arraystretch}{1.2}
    \begin{tabular}{>{\raggedright\arraybackslash}p{2.8cm} >{\raggedright\arraybackslash}p{7.2cm}}
        \toprule
        \textbf{Step} & \textbf{Typical Function} \\
        \midrule
        Perceive & Detect lanes, signs, drivable area, and dynamic obstacles. \\
        Localize & Estimate ego pose relative to map and road geometry. \\
        Predict & Estimate trajectories of nearby vehicles and pedestrians. \\
        Plan & Select route and short-horizon safe maneuver. \\
        Control & Convert to steering, throttle, and braking commands. \\
        Verify & Apply safety checks and fail-safe fallback if needed. \\
        \bottomrule
    \end{tabular}
\end{frame}

\begin{frame}{Case Insight from Core AI Texts}
    Introductory AI texts discuss robotic vehicles as complete AI agents:
    \begin{itemize}
        \item Driverless platforms combine \textbf{cameras, radar, and LiDAR}.
        \item Perception and planning operate continuously in closed loop.
        \item Sensor design strongly influences what can be perceived and controlled.
    \end{itemize}

    \vspace{0.2cm}
    This supports a key principle:
    \begin{quote}
        \small Better autonomy comes from strong integration of sensors, perception, planning, and control.
    \end{quote}
\end{frame}

\begin{frame}{Challenges in Real Deployment}
    \textbf{Technical challenges:}
    \begin{itemize}
        \item Long-tail edge cases are hard to cover in training data.
        \item Domain shift: weather, geography, and sensor aging.
        \item Real-time compute constraints on embedded hardware.
    \end{itemize}

    \vspace{0.2cm}
    \textbf{Safety and ethics challenges:}
    \begin{itemize}
        \item Reliable uncertainty estimation and fallback behavior.
        \item Bias and fairness in perception models.
        \item Accountability and regulation in autonomous decision making.
    \end{itemize}
\end{frame}

\begin{frame}{Evaluation Metrics (Overview)}
    Robotics and vision systems are evaluated at multiple levels:

    \vspace{0.1cm}
    \textbf{Perception:}
    \begin{itemize}
        \item Classification accuracy, precision, recall, F1.
        \item Detection mAP (mean Average Precision).
    \end{itemize}

    \vspace{0.1cm}
    \textbf{Autonomy/system:}
    \begin{itemize}
        \item Collision rate and safety intervention count.
        \item Lane-keeping and trajectory tracking error.
        \item Mission success and average completion time.
    \end{itemize}

    \vspace{0.1cm}
    \textit{The final objective is safe, robust behavior in real environments.}
\end{frame}

\begin{frame}{Summary and Next Steps}
    \textbf{What we covered:}
    \begin{itemize}
        \item Robotics as a sense-think-act AI system.
        \item Perception fundamentals and uncertainty.
        \item Sensor types and the need for fusion.
        \item Basic image classification and CNN intuition.
        \item Role of AI in autonomous systems and deployment risks.
    \end{itemize}

    \vspace{0.2cm}
    \textbf{Key takeaway:} Robust autonomy requires both strong models and strong systems engineering.

    \vspace{0.3cm}
    \textit{Use the next slides for quick self-assessment and discussion.}
\end{frame}

\begin{frame}{Self-Assessment Questions I}
    \begin{enumerate}
        \item Describe the \textit{sense-think-act loop} and explain why all three components are essential in robotics.
        \item Why is sensor fusion necessary? Give an example of how combining camera and LiDAR improves robustness.
        \item A CNN trained on sunny outdoor images shows poor performance on rainy images. What problem is this, and how might you address it?
    \end{enumerate}
\end{frame}

\begin{frame}{Self-Assessment Questions II}
    \begin{enumerate}
        \setcounter{enumi}{3}
        \item In autonomous driving, what trade-offs exist between classification confidence thresholds and system safety?
        \item Name at least two real-world challenges (not in training) that make robust autonomous systems hard to deploy.
    \end{enumerate}
\end{frame}

\end{document}
